% ==========================================================
% 080_role_in_frf.tex
% Forecast Admissibility Surface (FAS)
% ==========================================================

\section{Role of FAS in the Forecast Readiness Framework}
\label{sec:fas_role_frf}

The Forecast Admissibility Surface is a foundational governance construct within
the Forecast Readiness Framework (FRF). Its role is to determine, prior to
advanced modeling and evaluation, which slices of the forecast domain are
structurally admissible for participation in forecasting and readiness-control
workflows.

\subsection{Position in the FRF Lifecycle}

Within the FRF lifecycle, FAS is applied after the generation of a simple
baseline forecast and before the training or selection of higher-capacity model
classes. The baseline forecast provides the minimal signal required to expose
structural error anatomy, while avoiding the confounding effects of aggressive
fitting or feature interaction.

Once the admissibility surface is established, it is joined onto downstream
artifacts and used to constrain subsequent stages of the framework. Slices
classified as \Allowed\ proceed unimpeded through model training, evaluation, and
governance. Slices classified as \Conditional\ are subject to restricted
participation, while \Blocked\ slices are excluded entirely.

\subsection{Relationship to Other FRF Diagnostics}

FAS complements, but does not replace, other FRF diagnostics. Demand Quantization
Compatibility (DQC) assesses whether realized demand admits stable interpretation
at a given resolution. Forecast Primitive Compatibility (FPC) assesses whether a
forecast primitive can respond coherently to readiness intervention. Governance
policies and Readiness Adjustment Layers (RAL) determine how and when forecasts
are adjusted.

FAS operates orthogonally to these components. It does not evaluate measurement
validity (DQC), responsiveness (FPC), or policy outcomes. Instead, it answers a
more fundamental question: whether a slice should enter the forecasting and
governance stack at all. A slice may pass DQC and FPC checks yet still be excluded
by FAS due to structural hostility revealed by baseline error anatomy.

\subsection{Constraint Propagation}

Once established, the admissibility surface propagates constraints downstream in
a deterministic manner. Training pipelines may skip or downweight \Blocked\
slices. Evaluation metrics may exclude or flag \Conditional\ slices. Governance
decisions and readiness adjustments may be suppressed or capped based on
admissibility class.

Crucially, these constraints are not embedded implicitly within models or
metrics. They are enforced explicitly via the joined admissibility surface,
preserving transparency and interpretability across the FRF.

\subsection{Non-Substitutability}

FAS is not a substitute for careful modeling, evaluation, or governance. It does
not guarantee forecast quality, nor does it absolve downstream components from
their respective responsibilities. Its sole function is to prevent structurally
inadmissible slices from undermining the integrity of the forecasting system.

By explicitly positioning FAS within the FRF, the framework acknowledges a
critical operational reality: not all demand can or should be forecasted. Making
this boundary explicit strengthens the coherence, stability, and credibility of
the overall forecasting process.
